
The risk analysis section of the project plan is vital, identifying potential risks to mission success and outlining mitigation strategies. It summarizes major risks and steps taken to address them:

\begin{itemize}[leftmargin=1cm, itemindent=0.25cm, noitemsep, topsep=0pt, label=$\bullet$]
\item \textbf{Communication failure}: Critical for mission success, factors like interference, distance, and line of sight affect communication quality. Mitigate risks through system checks, tests, and deployment procedures, including onboard self-diagnostics and LED confirmation.
\item \textbf{Power limitations}: Payload power generation and storage are limited. Ensure power-efficient design and sufficient power for mission completion.
\item \textbf{Environmental factors}: Payload faces various environmental challenges like temperature, humidity, and vibration. Design for resilience and verify mission objectives through endurance, high G-force simulation, and landing tests.
\item \textbf{Technical issues}: Prototype development board to catch design mistakes early. Identify and prepare contingency plans for potential technical challenges.
\item \textbf{Budget constraints}: Ensure project feasibility within budget constraints, prioritizing mission objectives accordingly.
\end{itemize}

\subsection{Resource estimation}
To ensure that the project is completed successfully within the given budget and timeframe, it is crucial to perform a thorough estimation of the required resources. This includes identifying the necessary materials, equipment, and team members needed to complete the project, as well as the associated costs and time requirements. 

By performing a detailed resource estimation, the project team can effectively plan and allocate resources to ensure that the project is completed on time, within budget, and to the desired quality standards. 

Below is a breakdown of the resources required for the CanSat project:
\begin{itemize}[leftmargin=1cm, itemindent=0.25cm, noitemsep, topsep=0pt, label=$\bullet$]
\item Materials and components:
\begin{itemize}[label=\ding{111}, noitemsep, topsep=1pt]
\item Identify and list out all the required materials and components, including a microcontroller, sensors, batteries, antennas, radio modules, lightweight printable plastics, and other miscellaneous parts
\item Cost of materials can vary depending on the quality and quantity of components required
\end{itemize}
\item Tools and equipment:
\begin{itemize}[label=\ding{111}, noitemsep, topsep=1pt]
\item Various tools and equipment required, such as multimeters, 3d printers, oscilloscope, soldering station, and other miscellaneous parts
\end{itemize}
\item Software:
\begin{itemize}[label=\ding{111}, noitemsep, topsep=1pt]
\item Programming microcontrollers, designing 3D concepts, and analyzing data collected by the sensors
\item Using open-source software can significantly reduce the costs associated with software for the project.
\end{itemize}
\item Manpower and build time:
\begin{itemize}[label=\ding{111}, noitemsep, topsep=1pt]
\item Online tools and weekly meetings are used to ensure accuracy
\item Estimated effort of 630 hours for the entire project.
\end{itemize}
\item Testing and validation:
\begin{itemize}[label=\ding{111}, noitemsep, topsep=1pt]
\item Essential step to ensure the project performs as expected
\item May include the use of specialized equipment such as a drone or rocket
\end{itemize}
\item Shipping, transport, and logistics:
\begin{itemize}[label=\ding{111}, noitemsep, topsep=1pt]
\item Shipping and logistics costs may need to be factored into the project’s budget
\item May include shipping materials and components and travel costs for team members to attend the competition
\end{itemize}
\item Miscellaneous costs:
\begin{itemize}[label=\ding{111}, noitemsep, topsep=1pt]
\item Other costs that may need to be considered when estimating resources for the project
\end{itemize}
\end{itemize}

In conclusion, it is evident that a successful project outcome is heavily reliant on a proper estimation of resources required to complete the project. Neglecting the factors such as materials, components, tools, equipment, software, testing and validation, shipping and logistics, and other unforeseen costs can lead to budget overruns, delays, and ultimately failure to achieve project objectives.

\subsubsection{Budget}

We conducted thorough research to ensure that our budget aligns with the ideal specifications for our CanSat. This involved finding the most efficient and lightweight components without compromising performance, as well as maximizing battery life. By carefully considering each part, we aimed to optimize our budget and ensure that we could achieve our project objectives.

However, building the CanSat is just a part of the project. Therefore, our budget covers the entire project from start to finish, which can be divided into four main parts: 
\begin{itemize}[leftmargin=1cm, itemindent=0.25cm, noitemsep, topsep=0pt, label=$\bullet$]
    \item \textbf{Hardware parts}: all the necessary mechanical and electronic/electrical parts required to build both our custom development board and final payload. The cost estimation includes the purchase of a microcontroller, sensors, batteries, antennas, radio modules, lightweight printable plastics, and other miscellaneous parts.
    \item \textbf{Publicity and Branding}: covers the creation of custom t-shirts to promote our team and the event. In addition to custom t-shirts, we will also invest in creating promotional materials such as flyers, posters, and banners that showcase our team and the mission. These materials will be distributed both digitally and in person to help increase awareness of our project and the competition. Additionally, we will allocate a portion of our budget towards branding efforts such as website design and social media management to maintain a consistent and professional image for our team. Investing in Publicity and Branding will help us build a strong reputation within the community and attract potential sponsors and supporters to our cause.
    % \item \textbf{Travel Expenses}: includes the transportation costs, with tickets for the team to travel to the Competition Venue. This includes the cost of transportation by train, as well as any necessary local travel expenses.
    \item \textbf{Emergency Funds}: a contingency budget to cover any unforeseen costs that may arise during the project.
\end{itemize}

Below is a table that outlines all the anticipated expenses associated with the CanSat mission. This includes the costs of the components used to construct the CanSat, as well as any supplementary materials, equipment, and personnel necessary to complete the mission.

The components required for the project will be purchased from various suppliers such Optimus Digital, Cle\c{s}te.ro, Mouser, TME, Adafruit, and JLCPCB. It is worth noting that the team did not receive any kits from the organizers, and as a result, the team members did all the design and construction of the payload.

\begin{table}[htbp]
\centering
\arrayrulecolor{DeepSkyBlue4}
\begin{tabular}{>{\centering\arraybackslash}p{4.5cm}>{\centering\arraybackslash}lc}
\rowcolor{DeepSkyBlue4}
\hline
\multicolumn{1}{c}{\textbf{\color{white!50}{Category}}}& \multicolumn{1}{c}{\textbf{\color{white!50}{Cost Item}}} & \textbf{\color{white!50}{Cost \texteuro}} \\
\hline
\multicolumn{1}{l}{Hardware Parts (deductible) (134\texteuro)}& ESP32 ($\mu$C) & 6 \\
& \cellcolor{LightCyan1!50}LPS25H - Altimeter & \cellcolor{LightCyan1!50}25 \\
& MPU 9250 - IMU Sensor & 5 \\
& \cellcolor{LightCyan1!50}BME688 - Barometric Pressure & \cellcolor{LightCyan1!50}10 \\
& MCP9808 - Temperature Sensor & \cellcolor{LightCyan1!50}5 \\
& \cellcolor{LightCyan1!50}VEML6075 - UVA UVB UV Sensor & 5 \\
& Quectel L76 - GPS & 10 \\
& Power Switching regulators, LDO & 4 \\
& \cellcolor{LightCyan1!50}micro SD & \cellcolor{LightCyan1!50}6 \\
& Buzzer & 1 \\
& SMA Connectors & 2 \\
& \cellcolor{LightCyan1!50}Radio module & \cellcolor{LightCyan1!50}10 \\
& Parachute material & 20 \\
& Battery & 25 \\
\multicolumn{1}{l}{Consumables (non-deductible)}&  \cellcolor{LightCyan1!50}PCB components (RLC, diodes) & \cellcolor{LightCyan1!50}10 \\
& \cellcolor{LightCyan1!50}Parachute cords & \cellcolor{LightCyan1!50}5 \\
& \cellcolor{LightCyan1!50}Antennas & \cellcolor{LightCyan1!50}10 \\
& 3D printer filament & 15 \\
\multicolumn{1}{l}{Publicity \& Branding (100 \texteuro)}& \cellcolor{LightCyan1!50}T-Shirts & \cellcolor{LightCyan1!50}65 \\
& \cellcolor{LightCyan1!50}Flyers \& stickers & \cellcolor{LightCyan1!50}10 \\
% \multicolumn{1}{l}{Travel Expenses (315 \texteuro)} & Train tickets & 205 \\
% & \cellcolor{LightCyan1!50} Meals & \cellcolor{LightCyan1!50}110 \\
\multicolumn{1}{l}{Emergency Funds}& Miscellaneous & 100 \\
\hline
& \textbf{Total} & \textbf{349} \\
\hline
\end{tabular}
\caption{Foreseen costs for the CanSat mission.}
\label{tab:costs}
\end{table}




\subsubsection{External support}
The successful execution of the CanSat mission relies on the assistance and resources given by different organizations, departments, and companies. We are fortunate to have received sponsorship or in-kind support from the following entities:
\begin{itemize}[leftmargin=1cm, itemindent=0.25cm, noitemsep, topsep=0pt, label=$\bullet$]
    % \item \textbf{CoderDojo Oradea} has provided substantial financial aid, technical support, and feedback on both hardware and software. This makes them our biggest sponsor so far, allowing us to acquire all of the most expensive components of our CanSat and facilities where we can carry out our activities;
    \item \textbf{Funda\c{t}ia Comunitar\u{a} Oradea} has generously provided facilities where we can conduct our activities;
    % \item \textbf{EduTrust} has provided both financial support and working facilities for our team;
    % \item \textbf{Elektrobit} has provided financial support, access to facilities, equipment, and expertise in Embedded systems;
    \item \textbf{Depozitul de Tricouri} has agreed to provide customized T-shirts for each team member, which will help us promote our team and the event.
\end{itemize}


\subsection{Test plan}

The mission's success hinges on thorough pre-launch testing. This plan outlines tests to ensure the payload meets objectives and functions under various conditions, categorized into mechanical, software, communication, power, integration, ground station, battery charging, launch, and post-mission tests. The goal is to ensure accurate data collection and transmission, and survival during launch and landing. Annex G provides a testing roadmap with general ideas for each test.

Specific tests on electronics and structure are crucial to verify sensor functionality and structural integrity under launch, landing, and environmental conditions (see Annex H). Endurance tests assess the payload's ability to withstand harsh conditions, such as falls, vibration, shock, and prolonged G-forces during transport, providing insights into its performance limits (see Annex I).

By conducting these tests, the team ensures the CanSat is fully functional and mission-ready.

 


\subsection{Time management}
To keep the project on track and meet our timeline, we have divided the process into several phases, each with specific tasks and durations. We have completed the ideation phase and are currently in the design and prototyping stages. The upcoming phases include various testing campaigns and construction runs. The exact details of the testing phase will be determined as the project progresses, leading up to the final phase, which culminates in the competition.

We estimate a total of 630 man-hours needed for the entire project and have already invested 330 hours. The Gantt chart automatically updates as we advance through the project, allowing us to monitor progress and adjust our plans as necessary to ensure timely completion. By using this tool, we can stay on schedule and achieve our desired outcomes.

 

% We have started testing Jira in parallel with using TeamGantt.com to implement an Agile framework for future competitions. Jira is a project management software that provides a flexible and powerful platform for managing projects using an Agile framework. 

% By implementing Jira, we can create a central location for tracking tasks, issues, and progress throughout the project lifecycle. This will allow us to easily collaborate and communicate with team members and stakeholders, as well as prioritize tasks and adjust schedules based on changing requirements or circumstances. In addition, Jira offers a range of reporting and analytics features, providing valuable insights into team performance, task completion rates, and potential bottlenecks. With Jira, we can also track time and effort spent on each task, which will help us better estimate future efforts and improve our planning and scheduling for future competitions. By using both Jira and TeamGantt.com, we can optimize our project management process and ensure efficient, effective completion of the CanSat project.

To ensure that we complete all necessary tasks on time, the following is a summary of our timeline. For a more detailed timeline, please refer to the Gantt chart provided in the Appendix.

\begin{itemize}[leftmargin=1cm, itemindent=0.25cm, noitemsep, topsep=0pt, label=$\bullet$]
\item Team formation: 11.10.2023 - 19.11.2023
\begin{multicols}{2}[\vspace{-0.75\baselineskip}]
\begin{itemize}[label=\ding{59}, noitemsep, topsep=2pt]
\item Think of feasible mission objectives and requirements
% \item CanSat Application Work \& Registration
\end{itemize}
\end{multicols}
\vspace*{-0.75\baselineskip}
\item Ideation: 11.10.2023 - 28.04.2024
\begin{multicols}{2}[\vspace{-0.75\baselineskip}]
\begin{itemize}[label=\ding{59}, noitemsep, topsep=2pt]
\item Documentation writing (PDR, CDR, FDR)
\item Mission refinements
\end{itemize}
\end{multicols}
\vspace*{-0.75\baselineskip}
\item Design the Rocket and Payload: 03.01.2024 - 19.05.2024
\begin{multicols}{2}[\vspace{-0.75\baselineskip}]
\begin{itemize}[label=\ding{59}, noitemsep, topsep=2pt]
\item Mechanical structure 
\item Electrical system
\item Software
\item Recovery system
\item Ground support
\end{itemize}
\end{multicols}
\vspace*{-0.75\baselineskip}
\item Prototyping: 22.01.2024 - 17.04.2024
\item Testing (prototype): 01.04.2024 - 30.04.2024
\item Build the CanSat: 08.04.2024 - 15.05.2024
\item Launch Campaign: 29.05.2024 - 02.06.2024
\end{itemize}

% Our program takes into account the potential issues that may arise at each phase and the time required to address them. However, unforeseen events may occur during the project, and as such, we have developed measures for prevention and risk mitigation, which are detailed below.

% \begin{table}[htbp]
% \centering
% \arrayrulecolor{DeepSkyBlue4}
% \begin{tabular}{>{\raggedright\arraybackslash}p{8cm}>{\raggedright\arraybackslash}p{7cm}}
% \rowcolor{DeepSkyBlue4}
% \hline
% \multicolumn{1}{c}{\textbf{\color{white!50}{Major risks}}} & \multicolumn{1}{c}{\textbf{\color{white!50}{Mitigation}}} \\
% \hline
% Delays in obtaining mandatory materials and electronic components due to the global semiconductor chip shortage, resulting in unavailability and delivery delays & Place orders for materials and electronic components well in advance, find alternative components and adjust the project timeline accordingly \\
% \rowcolor{LightCyan1!50}Overtime required for the electrical design part due to thorough testing and prototyping of the onboard circuits before integration into the final design & Create a good electrical schematic and perform exhaustive checks before beginning the physical build to catch any errors \\
% Unforeseen technical difficulties that may require additional troubleshooting time and access to technical guidance from partners & Grant some extra time for troubleshooting and have access to technical guidance from our partners \\
% \rowcolor{LightCyan1!50}Adverse weather conditions during testing & Analyze the weather status in advance and plan at least one backup day for testing \\
% CanSat damage during testing due to inappropriate launch sites and lack of safety measures & Carefully choose the launch site when testing the CanSat and ensure that all safety measures are met \\
% \rowcolor{LightCyan1!50}Recovery failure due to poor GPS transmission and weak buzzer sound & Before launching, ensure that the GPS transmission is strong and the buzzer sound is loud \\
% Team member availability: Unexpected absences or departures of team members could delay the project and lead to workload imbalances & Clear communication and expectations among team members, cross-training team members on critical tasks, and having backup plans for key roles \\
% \hline
% \end{tabular}
% \caption{Unforeseen events and risk mitigation.}
% \label{tab:risks}
% \end{table}